%% LyX 2.4.4 created this file.  For more info, see https://www.lyx.org/.
%% Do not edit unless you really know what you are doing.
\documentclass[english]{article}
\usepackage[T1]{fontenc}
\usepackage[latin9]{inputenc}
\usepackage{enumitem}
\usepackage{amsmath}
\usepackage{amsthm}
\usepackage{cancel}
\usepackage{geometry}
\geometry{verbose,tmargin=1in,bmargin=1in,lmargin=1in,rmargin=1in}

\makeatletter

%%%%%%%%%%%%%%%%%%%%%%%%%%%%%% LyX specific LaTeX commands.
%% Because html converters don't know tabularnewline
\providecommand{\tabularnewline}{\\}

%%%%%%%%%%%%%%%%%%%%%%%%%%%%%% Textclass specific LaTeX commands.
\numberwithin{equation}{section}
\numberwithin{figure}{section}
\newlength{\lyxlabelwidth}      % auxiliary length 

%%%%%%%%%%%%%%%%%%%%%%%%%%%%%% User specified LaTeX commands.
\usepackage{hyperref}
\usepackage[usenames,dvipsnames,table]{xcolor} %adds additional color options
\hypersetup{colorlinks=true, urlcolor=black}

\usepackage{fancyhdr}
\usepackage{lastpage}
\pagestyle{fancy}
\usepackage{enumitem}
\usepackage{ifthen}
%sets math fonts to be more readable
\everymath=\expandafter{\the\everymath\displaystyle}
%\DeclareMathSizes{10}{10}{10}{10}
\usepackage{multicol}

\usepackage{tikz}
\usetikzlibrary{plotmarks}
\usepackage{pgfplots}
\pgfplotsset{compat=newest}

\usepackage{calc}
\usepackage[nomessages]{fp}

\usepackage{advdate}    % Advancing/saving dates
\usepackage{datetime}   % Dates formatting
\usepackage{datenumber} % Counters for dates

\usepackage{fancybox} %%ovalbox and fancy box settings
% Added by lyx2lyx
\setlength{\parskip}{\smallskipamount}
\setlength{\parindent}{0pt}

\makeatother

\usepackage{babel}
\begin{document}
\renewcommand{\author}{Dr.\,James Ashe}

\newcommand{\classNum}{107}

\newcommand{\classSec}{7 and 10}

\newcommand{\nameType}{}

\newcommand{\examType}{Exam 2 Review: sections 1.1-1.5}

\renewcommand{\date}{\hfill Exam 2 is on \formatdate{1}{10}{2014}}

%%First page's header%%%%%%%%%%%%%%%%%%%%%%
\fancypagestyle{firststyle} %%header settings for first pagr
{
	\fancyhead[L]{MTH \classNum{}(\classSec{}) \author{}\\
	\textbf{\examType{}} \date{}}
	\fancyhead[C]{\textbf{\nameType}}
	\setlength{\headsep}{14pt}
}
%%%%%%%%%%%%%%%%%%%%%%%%%%%%%%%%%%%%%%%%%%

%%Next page's header%%%%%%%%%%%%%%%%%%%%%%%
%\setlist{nolistsep}
\lhead{MTH \classNum{}(\classSec{})} 
\chead{\textbf{\nameType}} 
\cfoot{\thepage{}~of~\pageref{LastPage}} 
\setlength{\headsep}{5pt} 
%%%%%%%%%%%%%%%%%%%%%%%%%%%%%%%%%%%%%%%%%%%

%%Various settings; see the preamble for other settings%%%%%
%%\setenumerate[0]{leftmargin=12pt,itemindent=0pt,labelwidth=0pt}
%\setlist[2]{noitemsep}
%\setlist{nosep}
\setlist{itemsep=-3pt,topsep=-2pt}
\renewcommand{\labelenumi}{\textbf{\arabic{enumi}.}}
\renewcommand{\labelenumii}{\textbf{\alph{enumii}.}}
\thispagestyle{firststyle}
%%%%%%%%%%%%%%%%%%%%%%%%%%%%%%%%%%%%%%%%%%
\newcommand{\xmax}{10}
\newcommand{\xmin}{-10}
\newcommand{\xstep}{0} %must be xmin+n
\newcommand{\ymax}{10}
\newcommand{\ymin}{-10}
\newcommand{\ystep}{0} %must be ymin+n

\newcommand{\xspace}{\hspace*{40pt}}
\newcommand{\yspace}{\vspace*{.75cm}}
\newcommand{\rowColor}{gray!40}
\begin{enumerate}[resume]
\item Construct a Venn diagram to determine the validity of the given argument.

\begin{enumerate}
\item \vspace{-25pt}%
\begin{tabular}{l}
\tabularnewline
\tabularnewline
1. All basketball players are well conditioned.\tabularnewline
2. Bill is well-conditioned.\tabularnewline
\hline 
Therefore, Bill is a basketball player. \tabularnewline
\end{tabular}\yspace
\item \vspace{-25pt}%
\begin{tabular}{l}
\tabularnewline
\tabularnewline
1. All pilots must take a drug test.\tabularnewline
2. Some people who take drug tests get false positives.\tabularnewline
\hline 
Therefore, some pilots get false-positives. \tabularnewline
\end{tabular}\yspace
\item \vspace{-25pt}%
\begin{tabular}{l}
\tabularnewline
\tabularnewline
1. All squares are rectangles.\tabularnewline
2. Some quadrilaterals are squares.\tabularnewline
\hline 
Therefore, some quadrilaterals are rectangles. \tabularnewline
\end{tabular}\yspace
\end{enumerate}
\item Using the symbolic representations\\
\begin{tabular}{rl}
$b$: & I ride a bicycle to school.\tabularnewline
$s$: & It is snowing.\tabularnewline
$r$: & It is raining. \tabularnewline
\end{tabular}\\
express the following statements in symbolic form.

\begin{enumerate}
\item It is raining and I rode my bicycle to school.\yspace
\item If it is raining or snowing, then I will not ride my bicycle to school.\yspace
\end{enumerate}
\item Let $x$ be an integer. Using the symbolic representations\\
\begin{tabular}{rl}
$p$: & $x$ is a prime number.\tabularnewline
$o$: & $x$ is odd.\tabularnewline
$t$: & $x$ is equal to $2$.\tabularnewline
$e$: & $x$ is divisible by $8$.\tabularnewline
\end{tabular}\\
express the following in words.

\begin{enumerate}
\item $p\rightarrow(o\vee t)$\yspace
\item $e\rightarrow\sim p$\yspace
\item $\sim o\vee\sim e$\yspace
\item $\sim(o\wedge e)\wedge t$\yspace
\end{enumerate}
\item Construct a truth table for the symbolic expression.\begin{multicols}{2}

\begin{enumerate}
\item $p\wedge(q\rightarrow\sim p)$\\
\rowcolors{1}{white}{gray!50}%
\begin{tabular}{c|c|c|c|c}
$p$ & $q$ & \xspace & \xspace & \xspace\tabularnewline
\hline 
\hline 
T & T &  &  & \tabularnewline
\hline 
T & F &  &  & \tabularnewline
\hline 
F & T &  &  & \tabularnewline
\hline 
F & F &  &  & \tabularnewline
\end{tabular}\yspace
\item $(x\rightarrow y)\vee\sim x$\\
\rowcolors{1}{white}{gray!50}%
\begin{tabular}{c|c|c|c|c}
$x$ & $y$ & \xspace & \xspace & \xspace\tabularnewline
\hline 
\hline 
T & T &  &  & \tabularnewline
\hline 
T & F &  &  & \tabularnewline
\hline 
F & T &  &  & \tabularnewline
\hline 
F & F &  &  & \tabularnewline
\end{tabular}\yspace
\end{enumerate}
\end{multicols}
\begin{enumerate}[resume]
\item [\textbf{c.}]$\sim a\wedge(b\rightarrow c)$\\
\rowcolors{1}{white}{gray!50}%
\begin{tabular}{c|c|c|c|c|c}
$a$ & $b$ & \textbf{$c$} & \xspace & \xspace & \xspace\tabularnewline
\hline 
\hline 
T & T & T &  &  & \tabularnewline
\hline 
T & T & F &  &  & \tabularnewline
\hline 
T & F & T &  &  & \tabularnewline
\hline 
T & F & F &  &  & \tabularnewline
\hline 
F & T & T &  &  & \tabularnewline
\hline 
F & T & F &  &  & \tabularnewline
\hline 
F & F & T &  &  & \tabularnewline
\hline 
F & F & F &  &  & \tabularnewline
\end{tabular}\yspace
\end{enumerate}
\item Symbolically write the negation of the following statements.

\begin{enumerate}
\item The grass is not greener.\yspace\yspace
\item It is not the case that, if our car is not fixed then we cannot drive
to Asheville. \yspace\yspace
\end{enumerate}
\item Write the negation of the following statements.

\begin{enumerate}
\item Some freshmen are taking Math 107.\yspace\yspace
\item Some seniors are not taking Math 107.\yspace\yspace
\item All football players need to attend practice.\yspace\yspace
\item No professors attend football practice. \yspace\yspace
\end{enumerate}
\item Using the statement, ``If your order is over \$50, then you will
not have to pay shipping fees.'' write the following:

\begin{enumerate}
\item \textbf{Converse}:\yspace\yspace
\item \textbf{Inverse}:\yspace\yspace
\item \textbf{Contrapositive:}\yspace\yspace
\item \textbf{Contrapositive }in terms of a sufficient condition\textbf{:}\yspace\yspace
\item \textbf{Negation:}\yspace\yspace
\end{enumerate}
\item Using the statement, ``If the sun is going down, then the cows are
coming home.'' write the following:

\begin{enumerate}
\item \textbf{Converse}:\yspace\yspace
\item \textbf{Inverse}:\yspace\yspace
\item \textbf{Contrapositive:}\yspace\yspace
\item \textbf{Contrapositive }in terms of a necessary condition\textbf{:}\yspace\yspace
\item \textbf{Negation }using conjunction\textbf{:}\yspace\yspace
\end{enumerate}
\item Use a truth table to determine if the following arguments are valid.

\begin{enumerate}
\item \vspace{-25pt}%
\begin{tabular}{rll}
 &  & \tabularnewline
 &  & \tabularnewline
$b$: & You buy the bike. & 1. You buy the bike or the tent.\tabularnewline
$t$: & You buy the tent. & 2. You buy the bike.\tabularnewline
\cline{3-3}
 &  & Therefore, you do not buy the tent.\tabularnewline
\end{tabular}\\
\begin{tabular}{c|c|c|c|c|c|c}
\multicolumn{1}{c}{} & \multicolumn{1}{c}{} & \multicolumn{1}{c}{} & \multicolumn{1}{c}{} & \multicolumn{1}{c}{} & \multicolumn{1}{c}{} & \tabularnewline
\multicolumn{1}{c}{} & \multicolumn{1}{c}{} & \multicolumn{1}{c}{\ovalbox{$1$}} & \multicolumn{1}{c}{\ovalbox{$2$}} & \multicolumn{1}{c}{} & \multicolumn{1}{c}{\ovalbox{$c$}} & \ovalbox{Argument}\tabularnewline
$b$ & $t$ & \xspace & \xspace & \xspace & \xspace & \xspace\tabularnewline
\hline 
\hline 
\rowcolor{\rowColor}T & T &  &  &  &  & \tabularnewline
\hline 
T & F &  &  &  &  & \tabularnewline
\hline 
\rowcolor{\rowColor}F & T &  &  &  &  & \tabularnewline
\hline 
F & F &  &  &  &  & \tabularnewline
\end{tabular}\yspace\yspace
\item \vspace{-25pt} %
\begin{tabular}{rll}
 &  & \tabularnewline
 &  & \tabularnewline
$r$: & You register early. & 1. If you register early, then you will receive a discount.\tabularnewline
$d$: & You receive a discount. & 2. You register early.\tabularnewline
\cline{3-3}
 &  & Therefore, you will receive a discount.\tabularnewline
\end{tabular}\\
\begin{tabular}{c|c|c|c|c|c|c}
\multicolumn{1}{c}{} & \multicolumn{1}{c}{} & \multicolumn{1}{c}{} & \multicolumn{1}{c}{} & \multicolumn{1}{c}{} & \multicolumn{1}{c}{} & \tabularnewline
\multicolumn{1}{c}{} & \multicolumn{1}{c}{} & \multicolumn{1}{c}{\ovalbox{$1$}} & \multicolumn{1}{c}{\ovalbox{$2$}} & \multicolumn{1}{c}{} & \multicolumn{1}{c}{\ovalbox{$c$}} & \ovalbox{Argument}\tabularnewline
$r$ & $d$ & \xspace & \xspace & \xspace & \xspace & \xspace\tabularnewline
\hline 
\hline 
\rowcolor{\rowColor}T & T &  &  &  &  & \tabularnewline
\hline 
T & F &  &  &  &  & \tabularnewline
\hline 
\rowcolor{\rowColor}F & T &  &  &  &  & \tabularnewline
\hline 
F & F &  &  &  &  & \tabularnewline
\end{tabular}\yspace
\item \vspace{-25pt}%
\begin{tabular}{rll}
 &  & \tabularnewline
 &  & \tabularnewline
$b$: & You buy the property. & 1. If you buy the property, then you will need a loan.\tabularnewline
$n$: & You need a loan. & 2. If you need a loan, then you will pay 8\%.\tabularnewline
\cline{3-3}
$p$: & You will pay 8\%. & Therefore, if you buy the property then you will pay 8\%.\tabularnewline
\end{tabular}\\
\begin{tabular}{c|c|c|c|c|c|c|c}
\multicolumn{1}{c}{} & \multicolumn{1}{c}{} & \multicolumn{1}{c}{} & \multicolumn{1}{c}{} & \multicolumn{1}{c}{} & \multicolumn{1}{c}{} & \multicolumn{1}{c}{} & \tabularnewline
\multicolumn{1}{c}{} & \multicolumn{1}{c}{} & \multicolumn{1}{c}{} & \multicolumn{1}{c}{\ovalbox{$1$}} & \multicolumn{1}{c}{\ovalbox{$2$}} & \multicolumn{1}{c}{} & \multicolumn{1}{c}{\ovalbox{$c$}} & \ovalbox{Argument}\tabularnewline
$b$ & $n$ & $p$ & \xspace & \xspace & \xspace & \xspace & \xspace\tabularnewline
\hline 
\hline 
\rowcolor{\rowColor}T & T & T &  &  &  &  & \tabularnewline
\hline 
T & T & F &  &  &  &  & \tabularnewline
\hline 
\rowcolor{\rowColor}T & F & T &  &  &  &  & \tabularnewline
\hline 
T & F & F &  &  &  &  & \tabularnewline
\hline 
\rowcolor{\rowColor}F & T & T &  &  &  &  & \tabularnewline
\hline 
F & T & F &  &  &  &  & \tabularnewline
\hline 
\rowcolor{\rowColor}F & F & T &  &  &  &  & \tabularnewline
\hline 
F & F & F &  &  &  &  & \tabularnewline
\end{tabular}\yspace
\item \vspace{-25pt}%
\begin{tabular}{l}
\tabularnewline
\tabularnewline
1. Nothing intelligible ever puzzles me.\tabularnewline
2. Logic puzzles me.\tabularnewline
\hline 
Therefore, logic is unintelligible. \tabularnewline
\end{tabular}\\
\begin{tabular}{c|c|c|c|c|c|c|c|c|c}
\multicolumn{1}{c}{} & \multicolumn{1}{c}{} & \multicolumn{1}{c}{} & \multicolumn{1}{c}{} & \multicolumn{1}{c}{} & \multicolumn{1}{c}{} & \multicolumn{1}{c}{} & \multicolumn{1}{c}{} & \multicolumn{1}{c}{} & \tabularnewline
\multicolumn{1}{c}{} & \multicolumn{1}{c}{} & \multicolumn{1}{c}{} & \multicolumn{1}{c}{} & \multicolumn{1}{c}{} & \multicolumn{1}{c}{} & \multicolumn{1}{c}{} & \multicolumn{1}{c}{} & \multicolumn{1}{c}{} & \ovalbox{Argument}\tabularnewline
 &  &  & \xspace & \xspace & \xspace & \xspace & \xspace & \xspace & \xspace\tabularnewline
\hline 
\hline 
\rowcolor{\rowColor}T & T & T &  &  &  &  &  &  & \tabularnewline
\hline 
T & T & F &  &  &  &  &  &  & \tabularnewline
\hline 
\rowcolor{\rowColor}T & F & T &  &  &  &  &  &  & \tabularnewline
\hline 
T & F & F &  &  &  &  &  &  & \tabularnewline
\hline 
\rowcolor{\rowColor}F & T & T &  &  &  &  &  &  & \tabularnewline
\hline 
F & T & F &  &  &  &  &  &  & \tabularnewline
\hline 
\rowcolor{\rowColor}F & F & T &  &  &  &  &  &  & \tabularnewline
\hline 
F & F & F &  &  &  &  &  &  & \tabularnewline
\end{tabular}\yspace
\end{enumerate}
\end{enumerate}

\end{document}
